\section{What remains}\label{sec:outlook}
\paragraph{The gap of the jump form.} The second eigenvalue of $S_L$ sits near
$2\times10^{-4}$ and is roughly constant over $3/2\leq L\leq4$. That is a
spectral gap statement about an explicit jump Dirichlet form on an interval, it
is where the criticality of the target originates, and it is well posed. An
explanation of it would explain the collapse of the margin mechanically rather
than measuring it. The spectral theory of nonlocal Dirichlet forms is a
developed subject and has not been applied here.

\paragraph{The coincidence of Section~\ref{sec:necessary}.2.} The ground state
of an operator assembled from the archimedean density and the prime atoms is
$\cosh(x/2)$ to four digits, and $\cosh(x/2)$ is the pole direction. Nothing
proved here forces that. A proof would make the rank-one cancellation of the
pole term structural rather than numerical, and would very likely sharpen
Proposition~\ref{prop:necessary}.

\paragraph{The sharp form of the necessary condition.} $P_L$ has rank two, so
the exact criterion, rather than the relaxation \eqref{eq:necessary}, is a
condition on the two scalars $\ip{\cosh(x/2)}{S_L^{-1}\cosh(x/2)}$ and
$\ip{\sinh(x/2)}{S_L^{-1}\sinh(x/2)}$. It is exact and small, but $S_L^{-1}$ is
dominated by the near-null direction of Remark~\ref{rem:insufficient}, so the
two scalars are delicate in the same way the original problem is; this relocates
the difficulty rather than reducing it.

\paragraph{A cone adapted to the pole directions.} Section~\ref{sec:cone} closes
the pointwise cone and gives the reason. It does not close the possibility of a
cone adapted to the two pole directions, nor the complex-cone theory appropriate
to indefinite forms. The prize is named precisely: a Collatz--Wielandt pointwise
supersolution.

\paragraph{Scope.} None of the above is proposed as a route to the Riemann
hypothesis. The results of this paper are unconditional structure and one
necessary condition; what they are for is to constrain and to exclude candidate
sources cheaply, which is the use the companion investigation makes of them.

\subsection*{Preparation}
This working manuscript was produced with a large language model, Anthropic's
Claude (model \texttt{claude-opus-5}), working directly in the investigation
repository: the identities and proofs of
Sections~\ref{sec:density}--\ref{sec:necessary}, the check programs of
Appendix~\ref{app:checks}, and the drafting throughout. The research notes
accompanying the investigation name the model used for each, and record the
exploratory computations that are deliberately outside the registered checks.
Every statement here is classified in the text as a written proof, a one-sided
numerical value, or a labelled floating-point computation.
